\documentclass[12pt,a4paper]{article}

\usepackage[utf8]{inputenc}
\usepackage[T1]{fontenc}
\usepackage[polish]{babel}
\usepackage{lmodern}
\usepackage{amsmath,amssymb}
\usepackage{graphicx}
\usepackage{float}
\usepackage{geometry}
\usepackage{xcolor}
\usepackage{array}
\usepackage{hyperref}
\usepackage{listings}
\usepackage{tikz}
\usetikzlibrary{arrows.meta,positioning,shapes.geometric,fit}

\geometry{margin=2.5cm}
\raggedbottom
\frenchspacing
\clubpenalty=10000
\widowpenalty=10000
\displaywidowpenalty=10000
\emergencystretch=2em
\setlength{\textfloatsep}{0.75em plus 0.2em minus 0.2em}
\setlength{\floatsep}{0.65em plus 0.2em minus 0.2em}
\setlength{\intextsep}{0.65em plus 0.2em minus 0.2em}
\hypersetup{
    colorlinks=true,
    linkcolor=black,
    urlcolor=blue
}

\definecolor{noteline}{gray}{0.35}
\definecolor{codebg}{gray}{0.98}
\definecolor{codeframe}{gray}{0.35}
\definecolor{codecomment}{gray}{0.35}
\definecolor{graphfill}{gray}{0.96}
\definecolor{graphline}{gray}{0.30}

\lstdefinelanguage{Pseudo}{
    morekeywords={
        if,else,for,return,try,create,draw,keep,split,convert,
        update,solve,check,evaluate,improve,substitute,correct,
        remove,with,by,in,from,to,and,or
    },
    sensitive=false,
    morecomment=[l]{//},
    morestring=[b]"
}

\lstdefinestyle{pseudo}{
    language=Pseudo,
    basicstyle=\ttfamily\small,
    backgroundcolor=\color{codebg},
    frame=single,
    rulecolor=\color{codeframe},
    framerule=0.6pt,
    framesep=5pt,
    columns=fullflexible,
    keepspaces=true,
    breaklines=true,
    numbers=none,
    showstringspaces=false,
    keywordstyle=\bfseries,
    commentstyle=\itshape\color{codecomment},
    stringstyle=,
    captionpos=b,
    xleftmargin=0.5em,
    xrightmargin=0.5em,
    framexleftmargin=0.4em,
    framexrightmargin=0.4em
}

\renewcommand{\lstlistingname}{Algorytm}
\renewcommand{\lstlistlistingname}{Spis algorytmów}

\newcommand{\plainnote}[2]{%
\par\medskip
\noindent\begin{minipage}{0.94\linewidth}
\textbf{#1}\par\smallskip
#2
\end{minipage}
\par\medskip
}

\newcommand{\techbox}[2]{\plainnote{#1}{#2}}
\newcommand{\methodbox}[2]{\plainnote{#1}{#2}}
\newcommand{\warnbox}[2]{\plainnote{#1}{#2}}

\newcommand{\screenshot}[2]{%
\begin{figure}[H]
    \centering
    \IfFileExists{figures/#1}{%
        \includegraphics[width=0.82\linewidth,height=0.32\textheight,keepaspectratio]{figures/#1}%
    }{%
        \IfFileExists{docs/figures/#1}{%
            \includegraphics[width=0.82\linewidth,height=0.30\textheight,keepaspectratio]{docs/figures/#1}%
        }{%
            \fbox{%
                \begin{minipage}[c][4.2cm][c]{0.72\linewidth}
                    \centering
                    \textcolor{gray}{Ilustracja przykładowego wyniku: \texttt{\detokenize{#1}}}
                \end{minipage}
            }%
        }%
    }
    \caption{#2}
\end{figure}
}

\newcommand{\screenshotpair}[4]{%
\begin{figure}[H]
    \centering
    \begin{minipage}{0.47\linewidth}
        \centering
        \IfFileExists{figures/#1}{%
            \includegraphics[width=\linewidth,height=0.25\textheight,keepaspectratio]{figures/#1}%
        }{%
            \includegraphics[width=\linewidth,height=0.25\textheight,keepaspectratio]{docs/figures/#1}%
        }%
        \vspace{0.2em}

        {\small #2}
    \end{minipage}\hfill
    \begin{minipage}{0.47\linewidth}
        \centering
        \IfFileExists{figures/#3}{%
            \includegraphics[width=\linewidth,height=0.25\textheight,keepaspectratio]{figures/#3}%
        }{%
            \includegraphics[width=\linewidth,height=0.25\textheight,keepaspectratio]{docs/figures/#3}%
        }%
        \vspace{0.2em}

        {\small #4}
    \end{minipage}
    \caption{Dwa przykładowe wyniki działania modułu.}
\end{figure}
}

\title{\textbf{Dokumentacja projektu}\\Topological Visual Toolkit}
\author{Julia Bugajska \and Filip Jastrzębski \and Krzysztof Korzeniewski}
\date{}

\begin{document}

\maketitle

\begin{center}
    Repozytorium projektu: \url{https://github.com/DevFifi/Topological-Visual-Toolkit/}\\
    Działająca aplikacja: \url{https://topological-visual-toolkit.streamlit.app}
\end{center}

\section{Cel projektu}

Celem projektu było przygotowanie aplikacji z wygodnym interfejsem, która realizuje sześć zadań z przedmiotu \emph{Elementy topologii stosowanej}. Program został napisany jako aplikacja Streamlit. Poszczególne moduły pozwalają liczyć odległości w przestrzeniach metrycznych, odległości Czebyszewa funkcji, wielomiany Bernsteina oraz przybliżać obrazy i przeciwobrazy zbiorów.

W projekcie przyjęto podejście praktyczne: tam, gdzie da się otrzymać wynik symboliczny lub dokładny zapis, program próbuje go pokazać. W przypadkach wymagających metod numerycznych wynik jest opisany jako przybliżony. Jest to ważne szczególnie przy supremach i rysowaniu przeciwobrazów, ponieważ dla ogólnych wzorów często trzeba korzystać z metod przybliżonych.

\section{Ogólne działanie programu}

Aplikacja ma sześć widocznych podprogramów, wybieranych z panelu bocznego:

\begin{enumerate}
    \item \emph{Skończone przestrzenie metryczne},
    \item \emph{Odległość supremum na przedziale},
    \item \emph{Odległość supremum na prostokącie},
    \item \emph{Aproksymacja Bernsteina},
    \item \emph{Przeciwobraz funkcji skalarnej},
    \item \emph{Odwzorowania wektorowe}.
\end{enumerate}

\methodbox{Zgodność z treścią zadań}{
\small
\begin{tabular}{p{0.13\linewidth}p{0.76\linewidth}}
\textbf{Zadanie} & \textbf{Realizacja w aplikacji} \\
22 & \textbf{Skończone przestrzenie metryczne}: macierz odległości, średnica i odległość między zbiorami w $\mathbb{R}^n$. \\
19 & \textbf{Supremum na przedziale}: $d_\infty(f,g)$ na $[a,b]$ oraz punkty osiągania maksimum. \\
20 & \textbf{Supremum na prostokącie}: $d_\infty(f,g)$ na prostokącie w $\mathbb{R}^2$. \\
20 & \textbf{Aproksymacja Bernsteina}: wykres $f$, wykres $B_n(f)$, błąd i~animacja. \\
5 & \textbf{Funkcja $\mathbb{R}^2\to\mathbb{R}$}: sprawdzenie punktu i rysunek $f^{-1}(A)$. \\
6 & \textbf{Odwzorowania wektorowe}: rysunki $\Phi(C)$ oraz $\Phi^{-1}(B)$. \\
\end{tabular}
}

We wszystkich modułach używany jest wspólny sposób wpisywania wzorów. Użytkownik może pisać zarówno zwykłe wyrażenia, np. \texttt{sin(x)+x\string^2}, jak i~zapis podobny do~LaTeX-a, np. \texttt{\textbackslash frac\{1\}\{2\}}, \texttt{\textbackslash sqrt\{x\}}, \texttt{\textbackslash pi}, \texttt{e\string^x}. Po rozpoznaniu wzoru program pokazuje podgląd w~postaci matematycznej. W~panelu bocznym znajduje się też klawiatura matematyczna, która wstawia typowe symbole i~operatory do aktywnego pola tekstowego.

Program ma także lokalną pamięć przykładów. Każdy moduł korzysta ze swojej historii, dzięki czemu przykłady z~jednego zadania nie mieszają się z~przykładami z~innego. Wyjątkiem są pola tego samego typu w~ramach jednego modułu, np. zbiory $E$ i~$F$ w~module metrycznym mogą korzystać z~tej samej listy zapisanych zbiorów punktów.

\section{Wspólne techniki}

Do obliczeń symbolicznych i~numerycznych używane są biblioteki Pythona: przede wszystkim SymPy, NumPy, SciPy oraz Plotly. SymPy służy do~parsowania i~upraszczania wzorów, NumPy do~szybkich obliczeń tablicowych, SciPy do~optymalizacji numerycznej, a~Plotly do~interaktywnych wykresów. Funkcje wpisane przez użytkownika są zamieniane na~funkcje numeryczne przez \texttt{lambdify}; wartości nieskończone i~nieokreślone są odrzucane, żeby pojedynczy błędny punkt siatki nie psuł całego wyniku.

Parser wzorów obsługuje podstawowe działania, nawiasy, potęgowanie, pierwiastki, ułamki, stałe $e$ i $\pi$, funkcje trygonometryczne, funkcje wykładnicze i logarytmiczne. Parser zbiorów rozpoznaje między innymi przedziały, zbiory skończone, prostokąty, relacje oraz operacje logiczne sumy i przecięcia.

\methodbox{Mapa użytych technik}{
\small
\begin{tabular}{p{0.22\linewidth}p{0.66\linewidth}}
\textbf{Część programu} & \textbf{Główne techniki} \\
Parsery & normalizacja zapisu, rekurencyjne dzielenie po operatorach najwyższego poziomu, SymPy \\
Metryki & rachunek symboliczny, blokowe obliczenia NumPy, wzory dla pudełek i przedziałów \\
Suprema & siatka punktów, lokalna optymalizacja SciPy, osobne badanie brzegów \\
Bernstein & stabilna ewaluacja przez logarytmy, próbkowanie $f(k/n)$, numeryczny błąd supremum \\
Rysunki zbiorów & maski logiczne na siatce, kontury Plotly, linie ciągłe i~przerywane dla brzegu \\
\end{tabular}
}

\subsection{Obsługiwane rodzaje wejść}

W polach funkcji można wpisywać na przykład:
\begin{quote}
\small
\verb|x^2+y^2|,\quad
\verb|\sqrt{x}|,\quad
\verb|e^x|,\quad
\verb|\sin(\pi*x)|,\quad
\verb|\frac{1+x}{2}|.
\end{quote}
W module odwzorowań wektorowych osobno wpisuje się dwie współrzędne odwzorowania: $\Phi_1(x,y)$ i $\Phi_2(x,y)$.

Punkty i zbiory skończone można podawać w postaci list:
\begin{quote}
\small
\verb|(0,0), (1,1)|,\quad
\verb|[(0,0), (e,\sqrt{2})]|,\quad
\verb|{(0,0), (\pi,1/2)}|.
\end{quote}
W module metrycznym program sprawdza wymiar punktów i w razie zmiany $n$ dopasowuje liczbę współrzędnych.

Zbiory w $\mathbb{R}$ mogą być przedziałami i zbiorami skończonymi:
\begin{quote}
\small
\verb|[0,1]|,\quad
\verb|(0,1]|,\quad
\verb|(-oo,0)|,\quad
\verb|{0,1,\pi}|.
\end{quote}
Można je również łączyć przez sumę i przecięcie, np.
\begin{quote}
\small
\verb|[0,1] \cup {2}|,\quad
\verb|[0,2] \cap (1,3)|.
\end{quote}

Zbiory w $\mathbb{R}^2$ mogą być prostokątami, relacjami albo zbiorami punktów. Przykładowe poprawne zapisy to:
\begin{quote}
\small
\verb|[-1,1]x[-2,2]|,\quad
\verb|[-1,1]\times[-2,2]|,\\
\verb|x^2+y^2<=1|,\quad
\verb|1/4 < x^2+y^2 <= 1|,\\
\verb|{(0,0),(e,\sqrt{2})}|.
\end{quote}
Warunki można łączyć operatorami \verb|\cap|, \verb|\cup|, \verb|\land|, \verb|\lor|, a także zapisami \verb|and| i \verb|or|. W zbiorach docelowych dla odwzorowania $\Phi$ można używać zmiennych $u,v$ zamiast $x,y$.

W module metrycznym, oprócz skończonych list punktów, obsługiwane są też proste zbiory pudełkowe i ich kombinacje:
\begin{quote}
\small
\verb|[-1,1]x[0,2]|,\quad
\verb|[-1,1]x{0,1}x[0,2]|.
\end{quote}
Można używać także przecięć i sum takich składników.

Przy rysowaniu zbiorów w~$\mathbb{R}^2$ stosowana jest siatka punktów oraz dodatkowe wyznaczanie brzegu przez kontury funkcji. Dzięki temu kształty opisane nierównościami, np. krzywe algebraiczne, są wyświetlane znacznie dokładniej niż przy samym kolorowaniu pojedynczych pikseli siatki. Brzegi zbiorów otwartych są oznaczane linią przerywaną, a~domkniętych linią ciągłą, o~ile program jest w~stanie to rozpoznać z~wpisanego warunku.

\subsection{Architektura rozwiązania}

Kod został podzielony na trzy główne części. Warstwa \texttt{ui} odpowiada za formularze, podglądy LaTeX, historię wejść i wykresy. Warstwa \texttt{core} zawiera parsery, formatowanie wyników, bezpieczną ewaluację numeryczną i zapis stanu aplikacji. Warstwa \texttt{math\_modules} zawiera właściwe algorytmy matematyczne używane przez sześć podprogramów.

\techbox{Przepływ danych w typowym module}{
Użytkownik wpisuje wzór albo zbiór. Tekst jest normalizowany i parsowany. Jeżeli zapis jest poprawny, program pokazuje podgląd matematyczny. Po kliknięciu przycisku obliczeń funkcje są zamieniane na bezpieczne funkcje numeryczne, wykonywany jest algorytm modułu, a wynik trafia do wspólnego formatu zawierającego część dokładną, numeryczną, opis metody i ewentualne uwagi.
}

\begin{figure}[!htbp]
\centering
\begin{tikzpicture}[
    node distance=7mm and 7mm,
    >=Latex,
    every node/.style={font=\small},
    flowbox/.style={rectangle, rounded corners=1.5mm, draw=graphline, fill=graphfill, thick, align=center, minimum width=3.0cm, minimum height=0.9cm},
    calc/.style={rectangle, rounded corners=1.5mm, draw=graphline, fill=white, thick, align=center, minimum width=3.0cm, minimum height=0.9cm},
    store/.style={rectangle, rounded corners=1.5mm, draw=graphline, fill=white, thick, align=center, minimum width=2.8cm, minimum height=0.8cm}
]
\node[flowbox] (ui) {interfejs\\Streamlit};
\node[flowbox, right=of ui] (input) {pole wejścia\\i historia};
\node[flowbox, right=of input] (parser) {parser wzoru\\lub zbioru};
\node[calc, below=of parser] (module) {moduł\\matematyczny};
\node[calc, left=of module] (result) {\texttt{DualValue}\\i dane wykresu};
\node[flowbox, left=of result] (output) {LaTeX, liczba,\\notatki, wykres};
\node[store, above=of input] (state) {\texttt{app\_state.json}};

\draw[->, thick] (ui) -- (input);
\draw[->, thick] (input) -- (parser);
\draw[->, thick] (parser) -- (module);
\draw[->, thick] (module) -- (result);
\draw[->, thick] (result) -- (output);
\draw[<->, thick, dashed] (input) -- (state);
\end{tikzpicture}
\caption{Uproszczony przepływ danych w aplikacji.}
\end{figure}

\begin{lstlisting}[style=pseudo,caption={Ogólny schemat działania pól matematycznych}]
input_text
  -> normalizacja zapisu LaTeX-like
  -> parse_expr z dozwolonymi symbolami i funkcjami
  -> podglad LaTeX albo komunikat bledu
  -> obliczenia symboliczne i/lub numeryczne
  -> DualValue: exact, exact_latex, numeric, method, notes
  -> wynik w interfejsie
\end{lstlisting}

\subsection{Parser wyrażeń}

Parser wyrażeń przyjmuje praktyczny podzbiór zwykłego zapisu matematycznego i~zapisu podobnego do~LaTeX-a. Przed przekazaniem tekstu do~SymPy wykonywana jest normalizacja: potęgowanie \verb|^| jest traktowane jak potęgowanie, \verb|\frac{a}{b}| jest zamieniane na~iloraz, \verb|\sqrt{x}| na~pierwiastek, \verb|\pi| na~$\pi$, a~zapis \verb|e| na~stałą Eulera. Obsługiwane są też nawiasy \verb|\left|, \verb|\right|, funkcje trygonometryczne, logarytmy, wartość bezwzględna zapisana jako \verb+|x-y|+ oraz mnożenie domyślne, np. \verb|2x|.

\begin{lstlisting}[style=pseudo,caption={Normalizacja i parsowanie wyrażenia}]
parse_expression(text):
    if text is empty:
        return error
    text = replace_unicode_symbols(text)
    text = replace_latex_frac(text)
    text = replace_latex_sqrt(text)
    text = replace_latex_functions(text)
    text = replace_abs_bars(text)
    namespace = {x,y,u,v,x1..x200,y1..y200,e,pi,sin,cos,...}
    return sympy.parse_expr(text, implicit_multiplication, convert_xor)
\end{lstlisting}

\subsection{Parser zbiorów}

Parser zbiorów działa rekurencyjnie. Najpierw normalizuje operatory, np. \verb|\cap|, \verb|\land|, \verb|and| do przecięcia oraz \verb|\cup|, \verb|\lor|, \verb|or| do sumy. Następnie dzieli zapis tylko po operatorach znajdujących się na najwyższym poziomie nawiasów. Dzięki temu wyrażenie z przecinkiem w punkcie albo przedziale nie jest błędnie rozdzielane.

\begin{figure}[!htbp]
\centering
\begin{tikzpicture}[
    node distance=6mm and 9mm,
    >=Latex,
    every node/.style={font=\small},
    decision/.style={diamond, aspect=2.2, draw=graphline, fill=graphfill, thick, align=center, inner sep=1.5pt},
    action/.style={rectangle, rounded corners=1.5mm, draw=graphline, fill=white, thick, align=center, minimum width=3.2cm, minimum height=0.8cm},
    atom/.style={rectangle, rounded corners=1.5mm, draw=graphline, fill=white, thick, align=center, minimum width=3.2cm, minimum height=0.8cm}
]
\node[action] (norm) {normalizacja zapisu};
\node[decision, below=of norm] (or) {operator\\sumy?};
\node[action, right=of or] (union) {parsuj składniki\\i zbuduj sumę};
\node[decision, below=of or] (and) {operator\\przecięcia?};
\node[action, right=of and] (inter) {parsuj składniki\\i zbuduj przecięcie};
\node[decision, below=of and] (comma) {przecinek\\na zewnątrz?};
\node[action, right=of comma] (commaact) {traktuj jako\\przecięcie};
\node[atom, below=of comma] (atom) {przedział, punkty,\\prostokąt lub relacja};

\draw[->, thick] (norm) -- (or);
\draw[->, thick] (or) -- node[above] {tak} (union);
\draw[->, thick] (or) -- node[left] {nie} (and);
\draw[->, thick] (and) -- node[above] {tak} (inter);
\draw[->, thick] (and) -- node[left] {nie} (comma);
\draw[->, thick] (comma) -- node[above] {tak} (commaact);
\draw[->, thick] (comma) -- node[left] {nie} (atom);
\end{tikzpicture}
\caption{Schemat rekurencyjnego parsera zbiorów.}
\end{figure}

\begin{lstlisting}[style=pseudo,caption={Rekurencyjne parsowanie zbioru}]
parse_set(text):
    text = normalize_set_operators(text)
    if top_level_OR_exists(text):
        return union(parse_set(left), parse_set(right))
    if top_level_AND_exists(text):
        return intersection(parse_set(left), parse_set(right))
    if top_level_commas_exist(text):
        return intersection(parse_set(part_1), ..., parse_set(part_k))
    return parse_atom(text)

parse_atom(text):
    try interval, finite set, rectangle, finite point set
    try chained relation:  a < expr <= b
    try single relation:   lhs <= rhs
\end{lstlisting}

\techbox{Dlaczego dzielenie po najwyższym poziomie jest ważne}{
Zapis \texttt{\{(0,0),(1,1)\}} zawiera przecinki, ale nie oznacza przecięcia warunków. Program śledzi nawiasy okrągłe, kwadratowe i klamrowe, więc rozdziela tylko te operatory, które naprawdę stoją między składnikami zbioru.
}

\subsection{Dokładność, format wyników i bezpieczeństwo obliczeń}

Wyniki są przechowywane w małej strukturze \texttt{DualValue}. Może ona zawierać dokładny zapis symboliczny, zapis LaTeX, wartość numeryczną, opis metody oraz notatki o dokładności. Dzięki temu każdy moduł może rozróżnić wynik dokładny od wyniku numerycznego, zamiast wyświetlać wszystkie liczby tak samo.

Dokładne wyrażenia są upraszczane kilkoma metodami SymPy: \texttt{simplify}, \texttt{radsimp}, \texttt{sqrtdenest}, rozwijanie pierwiastków i faktoryzacja. Program wybiera wariant o najczytelniejszym zapisie według prostej funkcji oceny, która karze między innymi niepotrzebne kwadraty nawiasów.

\begin{lstlisting}[style=pseudo,caption={Dobór czytelniejszej postaci dokładnej}]
simplify_exact(expr):
    candidates = [
        expr,
        simplify(expr),
        radsimp(expr),
        sqrtdenest(expr),
        expand_roots(expr),
        factor(expand(simplify(expr)))
    ]
    remove failed transformations
    return candidate with shortest readable representation
\end{lstlisting}

Funkcje numeryczne są tworzone przez \texttt{lambdify}, ale są opakowane w funkcje bezpieczne. Jeżeli w jakimś punkcie pojawi się dzielenie przez zero, logarytm z liczby niedodatniej albo przepełnienie, wynik w tym punkcie jest zamieniany na \texttt{NaN}. Algorytmy szukające maksimum albo rysujące zbiory ignorują takie punkty.

\subsection{Pamięć przykładów i stan aplikacji}

Historia wejść jest zapisywana w pliku JSON. Każdy moduł ma osobne kategorie historii, dlatego przykłady funkcji z modułu Bernsteina nie mieszają się z przykładami zbiorów z modułu odwzorowań. Przy polach tego samego typu w jednym module historia jest wspólna, np. zbiory $E$ i $F$ w module metrycznym korzystają z tej samej listy zapisanych zbiorów punktów.

Program ma też plik stanu początkowego z przykładami. Przy pierwszym uruchomieniu albo przy przywracaniu stanu aplikacji dane z tego pliku są wczytywane do bieżącej pamięci. Przy aktualizacjach przykłady są scalane po identyfikatorze i po surowej wartości wpisu, żeby nie tworzyć zbędnych duplikatów.

\clearpage
\section{Podprogram 1: skończone przestrzenie metryczne}

\subsection{Opis zadania}

Dane są: liczba $n\in\mathbb{N}\setminus\{0\}$, metryka $d$ w przestrzeni $\mathbb{R}^n$ oraz niepusty skończony zbiór
\[
E=\{p_1,\ldots,p_s\}\subseteq\mathbb{R}^n.
\]
Program wypisuje macierz odległości
\[
D(p_1,\ldots,p_s)=\bigl[d(p_i,p_j)\bigr]_{i,j=1}^{s}
\]
oraz średnicę zbioru
\[
\operatorname{diam}(E)=\max_{p,q\in E} d(p,q).
\]
Dla dwóch niepustych zbiorów skończonych $E,F\subseteq\mathbb{R}^n$ program oblicza także
\[
\operatorname{dist}(E,F)=\inf\{d(p,q):p\in E,\ q\in F\}.
\]
Ponieważ zbiory są skończone, infimum jest w tym przypadku minimum.

\subsection{Metoda}

\subsubsection*{Dostępne metryki}

W programie dostępne są między innymi następujące metryki:
\[
\delta(x,y)=
\begin{cases}
0, & x=y,\\
1, & x\neq y,
\end{cases}
\]
\[
d_H(x,y)=\#\{i:x_i\neq y_i\},
\]
\[
d_p(x,y)=
\begin{cases}
\left(\sum_{i=1}^{n}|x_i-y_i|^p\right)^{1/p}, & p\geq 1,\\[0.4em]
\sum_{i=1}^{n}|x_i-y_i|^p, & 0<p<1,
\end{cases}
\]
oraz
\[
d_\infty(x,y)=\max_{1\leq i\leq n}|x_i-y_i|.
\]

Można również wpisać własny wzór metryki, np. sumę po współrzędnych. Program sprawdza poprawność wymiarów punktów, niepustość zbiorów oraz przypadki, w których wzór daje wartości nienumeryczne.

\subsubsection*{Obliczanie wyników}

Dla zwykłych skończonych list punktów program tworzy symetryczną macierz odległości. Wartości na przekątnej są równe $0$, a~dla par $i<j$ odległość jest liczona tylko raz i~przepisywana do~komórki symetrycznej. Jeżeli współrzędne punktów są symboliczne, np. zawierają $\sqrt{2}$, $\pi$ albo $e$, program najpierw próbuje zachować dokładny zapis, a~dopiero potem oblicza wartość dziesiętną.

Średnica zbioru skończonego jest maksimum po wszystkich parach punktów z~$E$, a~odległość dwóch zbiorów skończonych jest minimum po parach $(p,q)\in E\times F$. Dla bardzo dużych zbiorów albo dużego wymiaru program przełącza się na szybszą ścieżkę numeryczną: punkty są zamieniane na tablice NumPy, a~odległości są liczone blokami. Dzięki temu nie trzeba trzymać w~pamięci całej ogromnej macierzy odległości naraz.

\begin{lstlisting}[style=pseudo,caption={Macierz odległości i średnica zbioru skończonego}]
distance_matrix(E, d):
    for i = 1..s:
        D[i,i] = 0
    for i = 1..s:
        for j = i+1..s:
            value = d(p_i, p_j)
            D[i,j] = value
            D[j,i] = value
    return D

diameter(E, d):
    best = 0
    for every unordered pair (p_i, p_j):
        best = max(best, d(p_i, p_j))
    return best
\end{lstlisting}

\begin{lstlisting}[style=pseudo,caption={Szybka ścieżka dla dużych zbiorów}]
large_set_distance(E, F, d):
    convert E and F to numeric arrays
    split arrays into blocks of about 160 points
    for every block pair:
        compute pairwise distances by NumPy broadcasting
        update current minimum or maximum
    return best value and realizing pair
\end{lstlisting}

Jeżeli użytkownik poda zbiór w postaci przedziałowej lub pudełkowej, program nie wypisuje macierzy punktowej. W takim przypadku średnica i odległość są liczone z końców odpowiednich przedziałów. Dla odległości między pudełkami używany jest odstęp domknięć przedziałów w każdej współrzędnej, a dla średnicy największa możliwa różnica końców. Przy otwartych brzegach wynik może być supremum albo infimum nieosiąganym przez konkretne punkty, więc program opisuje to w komunikacie.

\techbox{Zbiory pudełkowe w module metrycznym}{
Dla produktu przedziałów program sprowadza odległość do listy różnic w kolejnych współrzędnych. Dla $\operatorname{dist}(E,F)$ używany jest odstęp między przedziałami, a dla $\operatorname{diam}(E)$ największa różnica między ich końcami. Następnie ta lista różnic trafia do wybranej metryki: $d_1$, $d_2$, $d_\infty$, $d_p$, Hamminga albo dyskretnej.
}

\subsection{Przykład}

Dla $n=2$, metryki euklidesowej i zbioru
\[
E=\{(0,0),(1,0),(0,1),(1,1)\}
\]
macierz odległości ma na przekątnej zera, a największa odległość wynosi $\sqrt{2}$. Program pokazuje wartości w czytelnej postaci LaTeX, np. $\sqrt{2}$ zamiast długiego zapisu technicznego.

\screenshotpair
{01_metryki_macierz.png}{Macierz odległości i średnica zbioru.}
{01_metryki_zbiory.png}{Odległość między dwoma zbiorami punktów.}

\clearpage
\section{Podprogram 2: odległość supremum na przedziale}

\subsection{Opis zadania}

Dany jest ograniczony przedział domknięty
\[
J=[a,b]\subseteq\mathbb{R}
\]
oraz funkcje ciągłe $f,g:J\to\mathbb{R}$. Program oblicza odległość Czebyszewa
\[
d_\infty(f,g)=\sup_{x\in J}|f(x)-g(x)|.
\]
Dodatkowo rysowany jest wykres pokazujący punkty, w których wartość $|f(x)-g(x)|$ jest równa lub bardzo bliska otrzymanemu supremum.

\subsection{Metoda}

Program najpierw bada funkcję
\[
h(x)=|f(x)-g(x)|.
\]
Sprawdzane są końce przedziału oraz punkty znalezione metodami numerycznymi. Najpierw wartości funkcji są liczone na gęstej siatce, a następnie najlepsze kandydaty są poprawiane lokalnie metodą optymalizacji jednej zmiennej. Dla prostszych wzorów program dodatkowo rozważa pochodną funkcji $(f-g)^2$, szuka punktów krytycznych i porównuje je z końcami przedziału. Jeżeli ta część symboliczna powiedzie się i zgadza się z wynikiem numerycznym, wynik jest pokazany także dokładnie.

\begin{lstlisting}[style=pseudo,caption={Supremum na przedziale}]
supremum_interval(f, g, [a,b]):
    h(x) = abs(f(x) - g(x))
    evaluate h on a dense grid
    candidates = grid maxima + endpoints
    for the best local candidates:
        improve candidate by bounded scalar optimization
    numeric_result = best value found

    if expression is not too complicated:
        solve derivative of (f-g)^2
        check roots inside [a,b] and endpoints
        if symbolic value agrees with numeric_result:
            return exact and numeric result
    return numeric result
\end{lstlisting}

\techbox{Dlaczego używamy $(f-g)^2$?}{
Maksima funkcji $|f-g|$ są takie same jak maksima $(f-g)^2$, a pochodna kwadratu jest zwykle wygodniejsza symbolicznie niż pochodna wartości bezwzględnej.
}

Przyjęta klasa funkcji to funkcje opisane jawnym wzorem, które da się obliczać na~przedziale $J$ i~które poza ewentualnymi punktami wykluczonymi przez dziedzinę nie powodują wartości nieskończonych lub \texttt{NaN}. Użytkownik powinien dobrać przedział tak, aby wpisane funkcje rzeczywiście były na~nim ciągłe.

\subsection{Przykład}

Dla
\[
f(x)=\sin x,\qquad g(x)=0,\qquad J=[0,2\pi]
\]
otrzymujemy
\[
d_\infty(f,g)=1.
\]
Maksimum jest osiągane w punktach, w których $\sin x=\pm 1$.

\screenshot{02_supremum_przedzial.png}{Odległość Czebyszewa dwóch funkcji na przedziale.}

\clearpage
\section{Podprogram 3: odległość supremum na prostokącie}

\subsection{Opis zadania}

Dany jest prostokąt domknięty
\[
P=[a,b]\times[c,d]\subseteq\mathbb{R}^2
\]
oraz funkcje ciągłe $f,g:P\to\mathbb{R}$. Program wyznacza przybliżoną odległość
\[
d_\infty(f,g)=\sup_{(x,y)\in P}|f(x,y)-g(x,y)|.
\]

\subsection{Metoda}

Podobnie jak w przypadku jednowymiarowym rozważana jest funkcja
\[
h(x,y)=|f(x,y)-g(x,y)|.
\]
Program sprawdza narożniki, brzegi oraz punkty wewnątrz prostokąta. Wewnątrz prostokąta najpierw liczona jest siatka wartości, potem kilka najlepszych punktów startowych jest poprawianych metodą L-BFGS-B z ograniczeniami do danego prostokąta. Brzegi są traktowane osobno jako cztery zadania jednowymiarowe. Ostateczny wynik jest największą wartością znalezioną w tych krokach.

\begin{lstlisting}[style=pseudo,caption={Supremum na prostokącie}]
supremum_rectangle(f, g, P):
    h(x,y) = abs(f(x,y) - g(x,y))
    evaluate h on a rectangular grid
    starts = several best grid points + corners
    for start in starts:
        improve by L-BFGS-B with bounds of P

    for each of four edges of P:
        solve one-dimensional maximum on the edge

    also check four corners
    return the largest finite value found
\end{lstlisting}

\techbox{Rola brzegów}{
Sama optymalizacja z punktów wewnętrznych może przeoczyć maksimum leżące na brzegu. Dlatego program osobno maksymalizuje funkcję na bokach prostokąta i osobno sprawdza narożniki.
}

\subsection{Przykład}

Dla
\[
f(x,y)=x^2+y^2,\qquad g(x,y)=0,\qquad P=[-1,1]\times[-1,1]
\]
największa wartość występuje w narożnikach prostokąta i wynosi $2$.

\screenshot{03_supremum_prostokat.png}{Odległość Czebyszewa funkcji na prostokącie.}

\clearpage
\section{Podprogram 4: aproksymacja Bernsteina}

\subsection{Opis zadania}

Dla funkcji ciągłej $f:[0,1]\to\mathbb{R}$ program pokazuje, jak wielomiany Bernsteina przybliżają funkcję $f$. Wielomian stopnia $n$ ma postać
\[
B_n(f)(x)=\sum_{k=0}^{n} f\left(\frac{k}{n}\right)
\binom{n}{k}x^k(1-x)^{n-k}.
\]

Aplikacja rysuje wykresy $f$ oraz $B_n(f)$, oblicza przybliżoną odległość
\[
d_\infty(f,B_n(f))=\sup_{x\in[0,1]}|f(x)-B_n(f)(x)|
\]
oraz tworzy animację pokazującą zmianę wykresu $B_n(f)$ przy wzroście $n$.

\subsection{Metoda}

Dla małych stopni program może pokazać pełną postać symboliczną wielomianu. W kodzie jest to ograniczone do niewielkich $n$, ponieważ pełny wzór bardzo szybko staje się nieczytelny. Dla większych stopni program wyświetla tylko kilka pierwszych składników sumy i wielokropek.

Numerycznie wielomian jest liczony z wartości $f(k/n)$ dla $k=0,\ldots,n$. Żeby uniknąć przepełnień przy dużych współczynnikach dwumianowych, składniki bazowe są obliczane przez logarytmy i funkcję gamma. Błąd $d_\infty(f,B_n(f))$ jest liczony tak samo jak supremum na przedziale: przez siatkę punktów na $[0,1]$ i lokalne doszukiwanie największej wartości błędu.

\begin{lstlisting}[style=pseudo,caption={Numeryczna ewaluacja wielomianu Bernsteina}]
bernstein(f, n):
    samples[k] = f(k/n) for k = 0..n

    B(x):
        result = 0
        for k = 0..n:
            log_basis =
                log(n!) - log(k!) - log((n-k)!)
                + k*log(x) + (n-k)*log(1-x)
            basis = exp(log_basis)
            correct basis at x=0 and x=1
            result += samples[k] * basis
        return result
\end{lstlisting}

\begin{lstlisting}[style=pseudo,caption={Animacja Bernsteina dla dużego n}]
animation_degrees(target_n, max_frames):
    if target_n <= max_frames:
        return [1,2,...,target_n]
    else:
        return max_frames degrees evenly spaced from 1 to target_n
\end{lstlisting}

\techbox{Dlaczego nie rozwijamy symbolicznie dużego $B_n$?}{
Dla $n=10000$ pełna suma ma $10001$ składników, więc jej wypisanie nie pomaga w~zrozumieniu wyniku. Program pokazuje początek wzoru, ale wykres, błąd i~animację liczy numerycznie.
}

\subsection{Przykład}

Jako przykład funkcji o większych oscylacjach można wziąć
\[
f(x)=\sin(50x)
\]
i porównać ją z wielomianem Bernsteina dla dużego stopnia, np. $n=10000$. W takim przypadku nie ma sensu wypisywać pełnego wzoru symbolicznego, ale wykres i błąd można obliczać numerycznie. Animacja pokazuje wtedy wybrane stopnie od $1$ do $10000$, dzięki czemu widać, jak przybliżenie poprawia się wraz ze wzrostem $n$.

\screenshotpair
{04_bernstein_wykres.png}{Wykres funkcji i wielomianu Bernsteina.}
{04_bernstein_animacja.png}{Animacja kolejnych wielomianów.}

\clearpage
\section{Podprogram 5: przeciwobraz funkcji skalarnej}

\subsection{Opis zadania}

Dana jest funkcja
\[
f:\mathbb{R}^2\to\mathbb{R},
\]
zbiór $A\subseteq\mathbb{R}$ oraz punkt $(x_0,y_0)\in\mathbb{R}^2$. Program sprawdza, czy
\[
(x_0,y_0)\in f^{-1}(A),
\]
czyli czy
\[
f(x_0,y_0)\in A.
\]
Oprócz tego aplikacja rysuje przybliżony przeciwobraz
\[
f^{-1}(A)=\{(x,y)\in\mathbb{R}^2:f(x,y)\in A\}.
\]

\subsection{Metoda}

Zbiór $A$ można podać jako przedział, sumę przedziałów, przecięcie przedziałów albo zbiór skończony. Najpierw program oblicza wartość $f(x_0,y_0)$ w~badanym punkcie i~sprawdza, czy należy ona do~$A$. Jeżeli da się to zrobić symbolicznie, używany jest zapis dokładny; w~przeciwnym razie program przechodzi do~sprawdzenia numerycznego.

Rysunek jest wykonywany w wybranym oknie widoku. Program tworzy siatkę punktów, liczy na niej wartości $f(x,y)$ i koloruje te punkty, dla których $f(x,y)\in A$. Dla przedziałów rysowane są także poziomice odpowiadające końcom przedziału. Brzeg przedziału domkniętego jest linią ciągłą, a brzeg przedziału otwartego linią przerywaną. Dla zbiorów skończonych rysowane są poziomice $f(x,y)=a$ dla kolejnych elementów $a\in A$.

\begin{lstlisting}[style=pseudo,caption={Przynależność punktu i rysowanie przeciwobrazu}]
scalar_preimage(f, A, P=(x0,y0)):
    value = f(x0,y0)
    try symbolic membership: value in A
    if symbolic result is undecidable:
        classify value numerically

    create grid in chosen drawing window
    Z = f(X,Y)
    mask = membership of Z in A
    draw filled contour of mask
    draw boundary contours f(x,y)=endpoint(A)
    use dashed line for open endpoints
\end{lstlisting}

\techbox{Interpretacja rysunku}{
Kolor wypełnienia oznacza punkty, które spełniają warunek $f(x,y)\in A$ w wybranym oknie. Linie brzegowe są poziomicami funkcji $f$, więc pokazują, gdzie wartość funkcji trafia w granice zbioru $A$.
}

\subsection{Przykład}

Jeżeli
\[
f(x,y)=x^2+y^2,\qquad A=[0,1],
\]
to przeciwobrazem jest domknięty dysk jednostkowy
\[
f^{-1}(A)=\{(x,y):x^2+y^2\leq 1\}.
\]

\screenshot{05_przeciwobraz_skalarny.png}{Przeciwobraz zbioru dla funkcji $f:\mathbb{R}^2\to\mathbb{R}$.}

\clearpage
\section{Podprogram 6: odwzorowania wektorowe}

\subsection{Opis zadania}

Dane jest odwzorowanie
\[
\Phi:\mathbb{R}^2\to\mathbb{R}^2
\]
oraz zbiory $B,C\subseteq\mathbb{R}^2$. Program rysuje przybliżony przeciwobraz
\[
\Phi^{-1}(B)=\{(x,y):\Phi(x,y)\in B\}
\]
oraz obraz
\[
\Phi(C)=\{\Phi(x,y):(x,y)\in C\}.
\]

\subsection{Metoda}

Przeciwobraz jest liczony przez sprawdzanie warunku przynależności $\Phi(x,y)\in B$ na siatce punktów w wybranym oknie. Jeżeli $B$ jest opisany relacją, np. $u^2+v^2\leq 1$, program podstawia do tej relacji współrzędne odwzorowania $\Phi_1(x,y)$ i $\Phi_2(x,y)$, a następnie rysuje obszar oraz brzeg przez poziomicę odpowiedniej różnicy stron równania.

Obraz zbioru $C$ jest rysowany przez próbkowanie punktów należących do $C$ i nanoszenie wartości $\Phi(x,y)$ na płaszczyznę obrazu. Dla odwzorowania identycznościowego program korzysta z dokładniejszego trybu rysowania samego zbioru i jego brzegu, ponieważ wtedy obraz jest po prostu tym samym zbiorem. Dla ogólnego odwzorowania używane jest próbkowanie, dlatego jakość zależy od wybranej rozdzielczości.

Program obsługuje zbiory opisane nierównościami, prostokąty, relacje oraz proste operacje na zbiorach, takie jak suma i przecięcie. Przy przecięciach brzeg jednego warunku jest przycinany drugim warunkiem, dzięki czemu na rysunku nie pojawiają się zbędne fragmenty brzegu poza właściwym zbiorem. Dla nierówności ostrych brzeg jest rysowany linią przerywaną, a dla nieostrych linią ciągłą. Użytkownik może zwiększyć jakość rysowania, jeżeli zależy mu na dokładniejszym obrazie kosztem czasu obliczeń.

\begin{lstlisting}[style=pseudo,caption={Przeciwobraz w module odwzorowań wektorowych}]
vector_preimage(Phi, B, window):
    create grid (X,Y) in source window
    U = Phi_1(X,Y)
    V = Phi_2(X,Y)
    mask = membership of (U,V) in B
    draw filled contour of mask

    if B is relation lhs(u,v) <= rhs(u,v):
        substitute u=Phi_1(x,y), v=Phi_2(x,y)
        draw boundary as contour lhs-rhs = 0
\end{lstlisting}

\begin{lstlisting}[style=pseudo,caption={Obraz zbioru w module odwzorowań wektorowych}]
vector_image(Phi, C, window):
    create grid in source window
    mask = membership of (X,Y) in C
    U = Phi_1(X,Y)
    V = Phi_2(X,Y)
    keep only points where mask is true and U,V are finite
    if too many points:
        thin the sample
    draw points (U,V) in target plane
\end{lstlisting}

\techbox{Specjalny przypadek odwzorowania identycznościowego}{
Jeżeli $\Phi(x,y)=(x,y)$, obraz zbioru $C$ jest równy samemu $C$. Program wykrywa ten przypadek i zamiast chmury punktów rysuje wypełniony obszar oraz brzeg, co daje znacznie lepszą jakość dla zbiorów zadanych nierównościami.
}

\subsection{Przykład}

Przykładowo można użyć odwzorowania
\[
\Phi(x,y)=
\left(
x\cos(x^2+y^2)-y\sin(x^2+y^2),\,
x\sin(x^2+y^2)+y\cos(x^2+y^2)
\right),
\]
czyli obrotu punktu o kąt zależny od jego odległości od początku układu. Dla zbioru
\[
C=\{(x,y): |x|\leq 1.6,\ |y|\leq 1.6,\ \sin(4x)\sin(4y)\geq 0\}
\]
oraz
\[
B=\{(u,v):0.8\leq u^2+v^2\leq 2.8,\ u\geq 0\}
\]
program rysuje obraz $\Phi(C)$ oraz przeciwobraz $\Phi^{-1}(B)$.

Do prostego sprawdzenia jakości rysowania można też użyć odwzorowania identycznościowego $\Phi(x,y)=(x,y)$ oraz zbioru
\[
(x^2+y^2-1)^3-x^2y^3<0,
\]
który daje charakterystyczny kształt serca.

\screenshot{06_odwzorowania_wektorowe.png}{Obraz i przeciwobraz dla odwzorowania $\Phi:\mathbb{R}^2\to\mathbb{R}^2$.}

\clearpage
\section{Ograniczenia programu}

Program nie próbuje rozwiązywać symbolicznie wszystkich możliwych problemów z analizy i topologii. Założono, że użytkownik podaje funkcje i zbiory w rozsądnej postaci: bez niejednoznacznych zapisów, z odpowiednio dobranym oknem rysowania oraz z dziedziną zgodną z treścią zadania. 

Najważniejsze ograniczenia są następujące:

\begin{itemize}
    \item dokładne wyniki są pokazywane tam, gdzie SymPy jest w stanie je sensownie uprościć,
    \item supremum dla trudnych funkcji jest liczone numerycznie,
    \item rysunki przeciwobrazów i obrazów są przybliżeniami zależnymi od siatki oraz wybranego okna,
    \item bardzo skomplikowane wzory mogą wymagać mniejszej rozdzielczości rysowania albo dłuższego czasu obliczeń,
    \item program zakłada, że funkcje podane w zadaniach o supremum są ciągłe na wybranych dziedzinach.
\end{itemize}

\end{document}
